% !TEX root = ../main.tex
\section{Introduction}
\label{sec:introduction}

Service robots are expected to become part of everyday life, assisting people at
home and at work with ordinary manipulation tasks. Meeting that expectation
requires a robot to handle objects it has never seen, in cluttered and
unstructured surroundings \cite{Mohammed2022ClutterReview}. Bin picking is the
canonical instance of this challenge: a robot must retrieve target items from a
dense, disordered container, where objects rest against one another and few are
graspable as they lie. Traditional pipelines that assume known object models and
a clear grasp pose become brittle here, because the required information is often
occluded or simply unavailable. Learning-based methods, and deep Reinforcement
Learning (RL) in particular, have therefore become the dominant approach to
manipulation in clutter \cite{Mohammed2022ClutterReview}.

A recurring insight in this literature is that grasping alone is not enough. When
objects are packed together, the robot must first rearrange them, and
non-prehensile actions such as pushing create the space and the poses that make a
later grasp feasible. Push-grasp synergy, in which pushing and grasping are learned
jointly, improves picking success in dense clutter beyond what grasping achieves
on its own \cite{Zeng2018VPG, Mokhtar2022SelfSupervisedLF, Kasaei2024HarnessingTS}.
This positions pushing as a foundational pre-grasp skill for service robots: the
ability to move an object into a target pose is what turns an ungraspable
configuration into a graspable one. This thesis studies that skill in isolation,
framed as a goal-reaching problem in which the agent must drive the object to a
target pose. In doing so it works toward the same broad goal as the
clutter-manipulation literature it builds on, namely competent, general
manipulation in unstructured environments.

Progress on this goal is limited less by what RL can express than by what it costs
to train. On-policy algorithms such as Proximal Policy Optimization (PPO) are
stable and widely used \cite{Schulman2017ProximalPO}, but they discard each batch
of experience after a single update. Their sample requirement is therefore high,
and it grows with the difficulty of the reward signal. When the reward is sparse
or poorly shaped, most of the interaction budget yields little useful gradient,
and the number of parallel environments needed for a stable update becomes the
dominant practical constraint. Reducing this sample cost, and understanding which
methods reduce it, is therefore a central concern for learning manipulation skills
of this kind.

Two families of techniques aim to reduce this constraint, and they act on
different parts of the problem. The first is reward shaping, which adds an
auxiliary signal that directs the agent toward useful behaviour without changing
what the task ultimately rewards. Potential-Based Reward Shaping (PBRS) is the
principled form of this idea \cite{Ng1999PolicyInvariance}: its shaping term is
the difference of a state potential, which leaves the optimal policy unchanged.
The second family modifies the training process rather than the reward. Curriculum
learning orders tasks from easy to hard \cite{Narvekar2020}, and its automatic
variants derive that order from the agent's own competence. Asymmetric Self-Play
(ASP) is the most prominent automatic form
\cite{Sukhbaatar2018IntrinsicMotivation, OpenAI2021AsymmetricSF}: one agent
proposes goals while a second learns to reach them, so the goal distribution
adapts to the learner without human specification. Both families are well
motivated. Reward shaping, in the potential-based form used here, produces a
single-agent policy that solves the task, and serves as the anchor of this study.
ASP is the method under scrutiny: it has succeeded on goal-reaching problems
elsewhere, including robotic manipulation \cite{OpenAI2021AsymmetricSF} and
reversible control tasks \cite{Sukhbaatar2018IntrinsicMotivation}, yet on the task
studied here it underperforms the single-agent anchor. That contrast is the puzzle
this thesis addresses. A proven method underperforms a simpler one, and the aim is
to explain why, by identifying the conditions of this task that differ from the
settings in which the method succeeded.

The task studied here is planar pushing, the non-prehensile skill motivated above,
reduced to a single object and framed as a goal-reaching problem so that the
learning question can be isolated. A robot must slide the object across a surface
to a target pose. Pushing suits this study
for a specific mechanical reason. Under the quasi-static assumption, a pushed
object's motion follows its friction limit surface, and the resulting twist
couples translation and rotation \cite{Mason1986Pushing, Goyal1991}. Almost every
push therefore changes both the position and the orientation of the object. A goal
that fixes both a target position and orientation is thus multi-objective, and the
two objectives cannot be optimised independently. This coupling makes planar
pushing a compact and demanding evaluation task: the reward must give an accurate
gradient for two interacting quantities, and any curriculum or self-play mechanism
must satisfy both criteria at once.

The experiments use a simulated Universal Robots UR5e arm with a Robotiq gripper,
pushing a T-shaped block on a tabletop in NVIDIA Isaac Lab
\cite{Mittal2023IsaacLab}. Isaac Lab runs many environment instances in parallel
on a single Graphical Processing Unit (GPU), which lets every model train at a
fixed compute budget on equal terms. The policy does not command joint torques
directly; instead, it selects an object-relative push primitive with four
parameters: an approach offset, an approach angle, a push length, and a push
direction. A cuRobo Inverse Kinematics (IK) solver then executes the chosen push.
Goals lie in the planar rigid-body configuration space $\mathrm{SE}(2)$, and a
push succeeds only when the object reaches a small position and orientation
tolerance of the goal. Chapter~\ref{section:methods} gives the full state, action,
and reward definitions.

Against this background, the thesis poses three research questions, together with
one supporting question on reward efficiency.
\begin{itemize}
  \item \textbf{RQ1:} Does asymmetric self-play, as applied successfully in prior
        robotic manipulation, succeed in this contact-rich multi-objective
        goal-reaching task under a single-GPU budget?
  \item \textbf{RQ2:} Which difference between this task and the settings where
        self-play succeeded, among goal coupling, reward formulation, imitation
        signal, goal representation, and compute scale, accounts for the
        difference in outcome?
  \item \textbf{RQ3:} Where in the learning dynamics does self-play break down:
        the coupled success gate, a biased value function, a suppressed imitation
        signal, or demonstrations that cannot be reproduced under stochastic
        contact?
  \item \textbf{Supporting RQ:} Can potential-based reward shaping reach
        single-agent success with fewer parallel environments than a hand-tuned
        improvement reward?
\end{itemize}

The first question asks whether the method transfers to this regime at all,
comparing self-play against the single-agent anchor under one budget and one
success criterion. The second question decomposes the difference between this task
and the successful settings into five candidate axes, and tests each in turn. The
comparison between a rotationally symmetric disc and the T-block isolates the
coupled goal as a single variable. The third question examines the mechanism, using
per-axis competence, the value-function bias, the imitation signal, and a
replay diagnostic that measures whether the goal-proposing agent's own
demonstrations reach the goals it sets. The supporting question isolates the
contribution of reward shaping alone.

\subsection{Contributions}

This thesis makes three contributions, in order of priority. The first is the
mechanism behind the self-play result. The models learn the individual position
and orientation sub-skills, but do not compose them under the coupled success gate;
the coupled multi-objective gate, rather than the reward or the goal encoder, is
the dominant cause. Five measurements support this: the disc-against-T-block
contrast, a per-axis-against-combined diagnostic, a value-bias measurement, an
imitation clip-saturation measurement, and a replay diagnostic. The second
contribution is the regime contrast. It sets this task against the settings where
self-play succeeded, along five axes: goal coupling, reward formulation, imitation
signal, goal representation, and compute scale. A negative result here is thereby
explained as a difference in regime, not a refutation of prior work. The third
contribution is a reward-efficiency result. Potential-based reward shaping reaches
single-agent success with roughly four times fewer parallel environments than the
improvement reward, which anchors the task as solvable. The practical message
follows: principled reward design should come before changes to the training
dynamics. We scope every self-play result to this regime. Asymmetric self-play
succeeds in the settings prior work used; the contribution here is the analysis of
why it underperforms in this one.

\subsection{Thesis Outline}

The remainder of this thesis is organised as follows.
Chapter~\ref{sec:drl_robotics} reviews the background and related work, covering
RL and PPO, potential-based reward shaping, curriculum learning and self-play,
goal-conditioned learning, and the mechanics of pushing. That chapter also draws
the mechanical contrast between the disc and the T-block, and explains the
multi-objective nature of the goal. Chapter~\ref{section:methods} defines the task
and its Markov Decision Process, explains why a naive improvement reward
underperforms, and derives the potential-based reward and its $\mathrm{SE}(2)$
form. It then describes the nine models grouped by role, and sets out the
experiment set that produces the results. Chapter~\ref{sec:experiments} presents
the experimental setup: the simulation and training configuration, the phases of
the experiment set, the held-out validation protocol, the cross-environment
comparison, and an off-policy baseline. Chapter~\ref{sec:results} reports the
results for each research question, namely whether self-play succeeds here, which
regime axis accounts for the outcome, the underlying mechanism, and the
environment-efficiency of shaping. Chapter~\ref{sec:conclusion} discusses the
findings, states the limitations, and outlines future work. The appendices collect
training curves, parameter sensitivity, scene descriptions, and detailed
diagnostics.
